\section{Текущие результаты}

Реализованный процесс уже поддерживает сбор корпуса, генерацию векторных представлений, извлечение ключевых слов, построение запросы, поиск эталонных ответов, генерацию кандидатов и автоматическая попарная оценка. Оставшийся отсутствующий этап --- собственно обучения с подкреплением обновление цикл. Поэтому данный раздел является узким: он описывает завершенные артефакты данных и фиксирует конкретный экспериментальный постановка для первых MRVF обучающие запуски.

Результат на уровне корпуса --- объединенная коллекция шуток примерно из \(664\text{k}\) объектов. В финальной схема обучения этот полный корпус служит пул поиска: все шутки кодируются векторных представлений и индексируются, а эталонные ответы для слабые запросы извлекаются из этого полного пространства векторных представлений. Однако извлечение ключевых слов нужно применять только к достаточно большому подмножеству, чтобы достичь целевого объем выборки. В нашем случае \(3-5\text{k}\) группы ключевых слов могут дать \(10-20\text{k}\) эталонные ответы для обучение, валидация и тестирование.

Для предварительный этап этого достаточно. Важный результат состоит в том, что данные и экспериментальный постановка теперь финализированы: корпус существует, эталонный ответ коллекция реализована, методология зафиксирована, а оставшаяся работа ограничена цикл обучения MRVF.

\section{Планируемая процедура обучения}

Репозиторий уже реализует полный процесс подготовки данных на Python с использованием Hugging Face Наборы данных, Parquet артефакты, FAISS-based ближайших соседей поиск и OpenAI-compatible API вызовы для векторное представление генерация, генерацию кандидатов и автоматический оценка. Однако для собственно обучающие запуски стек программного обеспечения немного изменится.

Планируемая реализация будет использовать модифицированную версию \texttt{trl} для оптимизация и \texttt{vLLM} для генерация. \texttt{trl} предоставляет со стороны обучения логика для оптимизация стратегии, тогда как \texttt{vLLM} будет использоваться и для генерации ответы базовых моделей, и для обслуживания траектории генерации во время обучение с подкреплением. Поскольку MRVF зависит от выборка нескольких траектории рассуждения на запрос, пропускная способность генерации траекторий важен для эффективности.

Целевые модели относятся к семейству Qwen3 в диапазоне размеров \(1.7\text{B}\) и \(4\text{B}\). Для каждого размера планируются следующие базовые модели:
\begin{enumerate}[leftmargin=1.8em]
 \item \textbf{Qwen3 Instruct}, базовая модель в без режима рассуждения режим;
 \item \textbf{Qwen3 Thinking}, базовая модель в thinking режим;
 \item \textbf{MRVF-обученных Qwen3 Thinking}, предлагаемая модель;
\end{enumerate}

Такая постановка сохраняет интерпретируемость сравнения. Он изолирует три отдельных вопроса: достаточно ли обычного инструкция настройка для генерации юмора; улучшает ли родной thinking режим генерацию юмора уже без дополнительного обучения; и может ли MRVF дополнительно улучшить thinking модель.

Все обучающие запуски планируются на 1--2 \texttt{NVIDIA H200} GPUs в HSE HPC кластер. Этого достаточно для целевых размеров моделей и позволяет распараллеливать эксперименты между объекты.

Сравнение будет использовать процесс оценки, уже реализованный в репозитории. Ответы сопоставляются для одного запроса, порядок левого и правого вариантов рандомизируется, оценщик принуждается выбрать ровно одного победитель, а оценки моделей агрегируются через Bradley-Terry рейтинг. Этого достаточно для эксперимента, поскольку цель состоит в обнаружении того, дает ли MRVF измеримый сдвиг относительно instruct и базовые модели в режиме рассуждения. Более крупная многомерный оценка с оценивание по рубрике и человеческую разметку планируется после завершения обучения.

\section{Планируемые абляции}

План абляций напрямую связан с теорией. Наиболее информативные абляции следующие:
\begin{itemize}[leftmargin=1.5em]
 \item \textbf{Число наблюдаемых эталонных ответов на запрос \(|Y^{*}|\).} Сравнить меньшие и большие наборы эталонных ответов, например через обрезку результатов поиска до \(1\), \(3\) или \(5\) объектов.
 \item \textbf{коэффициент KL \(\beta\).} Проверить, насколько сильно модель должна быть удерживаться рядом с начальной контрольной точкой, чтобы компенсировать неполнота.
 \item \textbf{Размер группы \(G\).} Варьировать число выбранный траектории рассуждения на один запрос, чтобы оценить variance / пропускная способность компромисс MRVF оцениватель.
 \item \textbf{Размер модели.} Сравнить \(1.7\text{B}\) и \(4\text{B}\), а также проверить, насколько MRVF устойчив при применении к меньшим моделям.
\end{itemize}
